\raggedbottom

\makeatletter
\renewcommand\section{\@startsection {section}{1}{\z@}%
                                   {-4.2ex \@plus -1ex \@minus -.2ex}%
                                   {2.8ex \@plus.2ex}%
                                   {\normalfont\moderate\bfseries}}
\renewcommand\subsection{\@startsection{subsection}{2}{\z@}%
                                     {-3.8ex\@plus -1ex \@minus -.2ex}%
                                     {2.2ex \@plus .2ex}%
                                     {\normalfont\normalsize\bfseries\leftskip 40pt}}
\makeatother

\chapter{Project Defense and Project Expo}
\label{chap:project-defense-expo}

This chapter presents the evaluation and dissemination activities of
the Vid-Scratch project through two rounds of project defense and
participation in the SKE Project Expo. It summarizes the feedback
received from the evaluators, the actions taken by the team in response
to that feedback, and the outcomes of those improvements. In addition,
this chapter describes the project poster, booth presentation, visitor
questions, and reflections from presenting Vid-Scratch to a broader
audience.

Overall, the defense and expo activities helped the team clarify the
research direction of the project, improve the Software Design Document
(SDD), and communicate the concept of adversarial video poisoning more
effectively to both technical and non-technical audiences.


\vspace{1.2\baselineskip}
\section{First Project Defense}
\label{section:first-defense}

\subsection{Comments and Feedback from the Evaluators}
\label{subsection:first-defense-feedback}

During the first project defense, the evaluators provided relatively
few implementation-specific comments because Vid-Scratch is a
research-based project. Instead of focusing on user interface details
or ordinary application features, the main recommendation was for the
team to conduct deeper research before finalizing the technical
approach. The evaluators encouraged the team to study related work in
adversarial machine learning, video action recognition, and data
poisoning in order to identify a method that was both technically
sound and suitable for the project's scope.

The feedback emphasized that the success of the project depended on
selecting an appropriate research method rather than simply building a
desktop application. Since adversarial video poisoning is a specialized
topic, the team needed to understand the assumptions, limitations, and
experimental requirements of different attack strategies before
committing to an implementation.

\subsection{Revisions Required, Actions Taken, and Outcomes}
\label{subsection:first-defense-revisions}

In response to the first defense feedback, the team expanded the
research process and performed additional laboratory experiments. The
team reviewed methods related to adversarial perturbations, spatial
transform attacks, video classification models, and training-time
poisoning. The team also conducted AI laboratory work to build the
required background knowledge for implementing and evaluating the
pipeline.

As a result, the team gained a clearer understanding of the technical
problem and was able to select a more suitable method for Vid-Scratch.
The final direction focused on adversarial video poisoning for
temporal action recognition, using SlowFast R101 as the target model
and combining frame selection, adversarial generation, and evaluation
metrics such as SSIM and attack success rate. This outcome gave the
project a stronger research foundation and allowed the implementation
to proceed with a clearer technical scope.

\vspace{1.2\baselineskip}
\section{Second Project Defense}
\label{section:second-defense}

\subsection{Comments and Feedback from the Evaluators}
\label{subsection:second-defense-feedback}

In the second project defense, most of the feedback focused on the SDD
rather than the core technical content. The evaluators indicated that
some descriptions in the document were too strong or overstated when
compared with the actual implementation and experimental evidence. The
main concern was not that the project direction was incorrect, but
that the written explanation needed to be more precise, realistic, and
consistent across chapters.

This feedback differed from the first defense. In the first defense,
the main issue was the need for more research and experimentation. In
the second defense, the technical direction had already become clearer,
so the evaluators focused more on how accurately the project was
documented in the report.

\subsection{Additional Revisions, Actions Taken, and Final Outcomes}
\label{subsection:second-defense-revisions}

Following the second defense, the team revised the SDD across all
chapters. Statements that were too broad or exaggerated were rewritten
to better match the actual system behavior, implementation details,
and evaluation results. The document was adjusted to present
Vid-Scratch as an adversarial video poisoning pipeline with defined
experimental conditions, rather than making claims beyond the evidence
available from the project.

After the revisions, the updated SDD was reviewed again by the project
advisor. The final version passed the review, indicating that the
document more accurately represented the project scope, research
contribution, system design, and experimental findings.

\subsection{Suggestions on How Project Defense Can Better Support Software Development}
\label{subsection:defense-suggestions}

From the team's experience, the project defense process was useful for
identifying both research direction issues and documentation issues.
For research-based software projects, defense sessions can better
support development by encouraging teams to clearly connect their
implementation decisions with research evidence. This helps ensure that
the project is not only functional as software, but also properly
grounded in existing technical knowledge.

The defense process can also help teams improve the clarity and
accuracy of their reports. In the case of Vid-Scratch, the second
defense helped the team refine the SDD so that the claims, results,
and explanations were better aligned with the actual work completed.

\vspace{1.2\baselineskip}
\section{Project Expo Participation}
\label{section:project-expo}

\subsection{Poster Presentation and Poster Components}
\label{subsection:poster-presentation}

Vid-Scratch was presented at the SKE Project Expo using a project
poster and a live demonstration booth. The poster was designed to
communicate the core idea of the project clearly: Vid-Scratch is an
adversarial video poisoning pipeline for evaluating the robustness of
temporal action recognition models.

\begin{figure}[H]
    \centering
    \includegraphics[height=0.82\textheight,keepaspectratio]{chapter_8/project_poster.png}
    \caption{Vid-Scratch project poster presented at the SKE Project Expo.}
    \label{fig:vidscratch-poster}
\end{figure}

The poster included the background and motivation of the project,
explaining why video-based action recognition models may be vulnerable
to small adversarial perturbations. It also presented the key
contributions of Vid-Scratch, including an accessible desktop tool, a
Bayesian frame selection process, and a unified pipeline for
inference-time and training-time attack scenarios.

The poster further summarized the adversarial poisoning pipeline,
application interface, evaluation results, system architecture,
technology stack, and possible future work. These sections helped
visitors understand both the technical research problem and the
practical implementation of the system.

\subsection{Event Photos}
\label{subsection:event-photos}

The following photographs show the Vid-Scratch booth during the SKE
Project Expo. The booth included the project poster, a monitor for
demonstrating the application, and team members who explained the
system to visitors.

\begin{figure}[H]
    \centering
    \includegraphics[width=0.82\textwidth]{chapter_8/poster_photo.jpg}
    \caption{Team members presenting the Vid-Scratch poster.}
    \label{fig:poster-photo}
\end{figure}

\begin{figure}[H]
    \centering
    \includegraphics[width=0.9\textwidth]{chapter_8/booth_demo.jpg}
    \caption{Vid-Scratch booth setup with the application demonstration screen.}
    \label{fig:booth-demo}
\end{figure}

\begin{figure}[H]
    \centering
    \includegraphics[width=0.9\textwidth]{chapter_8/visitor_demo.jpg}
    \caption{Live demonstration and discussion with visitors during the expo.}
    \label{fig:visitor-demo}
\end{figure}

\begin{figure}[H]
    \centering
    \includegraphics[width=0.9\textwidth]{chapter_8/visitor_discussion.jpg}
    \caption{Explaining the Vid-Scratch poisoning workflow to expo visitors.}
    \label{fig:visitor-discussion}
\end{figure}

\subsection{Questions Raised During the Expo and the Team's Responses}
\label{subsection:expo-questions}

During the SKE Project Expo, most visitors asked what Vid-Scratch can
do and how the system works. The team explained that Vid-Scratch takes
an input video and generates a poisoned output video that appears
visually similar to humans, while causing a target action recognition
model to misclassify the action.

Visitors also asked about the full processing flow. In response, the
team explained the pipeline step by step, including video loading,
frame sampling, poison budget calculation, Bayesian frame selection,
adversarial generation, output mapping, and evaluation. Since many
visitors were unfamiliar with adversarial video poisoning, the team
focused on explaining why small perturbations can affect machine
learning models even when the visual difference is difficult for
humans to notice.

Another common question was the processing time required for each
video. The team answered by referring to the evaluation results and
explained that processing time depends on video length, hardware, and
pipeline settings. Based on the project benchmark, the average
processing time was approximately 42 seconds per video under the
tested configuration.

Overall, the team was able to answer the questions raised by visitors.
The discussion at the expo helped confirm that the project concept was
interesting, but also showed that adversarial video poisoning requires
clear explanation because it is not a familiar topic for most users.

\subsection{Recommendations for Improving the Expo}
\label{subsection:expo-recommendations}

The feedback from expo visitors was generally positive, and there were
no major suggestions for changing the current system. Some instructors
who visited the booth also commented that the project was already in a
good state. However, based on the team's own reflection, one useful
future improvement would be to expand the benchmark section.

A broader benchmark could compare Vid-Scratch against additional
videos, attack settings, hardware environments, or target video models.
This would make the evaluation stronger and help future audiences
understand the performance and limitations of the system more clearly.
In addition, benchmark results could make the project easier to explain
during public demonstrations because visitors could quickly see how
the system performs under different conditions.

Overall, the SKE Project Expo gave the team an opportunity to present
Vid-Scratch outside the defense setting, explain the research problem
to a wider audience, and validate that the project can be communicated
effectively through a combination of poster, live demonstration, and
step-by-step explanation.

\makeatletter
\renewcommand\section{\@startsection {section}{1}{\z@}%
                                   {-3.5ex \@plus -1ex \@minus -.2ex}%
                                   {2.3ex \@plus.2ex}%
                                   {\normalfont\moderate\bfseries}}
\renewcommand\subsection{\@startsection{subsection}{2}{\z@}%
                                     {-3.25ex\@plus -1ex \@minus -.2ex}%
                                     {1.5ex \@plus .2ex}%
                                     {\normalfont\normalsize\bfseries\leftskip 40pt}}
\makeatother

\flushbottom
